\uuid{laWJ}
\titre{Tangentes}
\auteur{hanani ; mzari}
\organisation{exo7}
\licence{CC BY-NC-SA 4.0}
\datecreate{2020-01-20}
\niveau{L1}
\module{Géométrie}
\chapitre{Courbes planes}
\sousChapitre{Courbes paramétrées}
\difficulte{2}
\theme{Courbes paramétrées, tangentes}
\exo7qcmid{lille-2.courbes.004}

\contenu{
\begin{question}
\qclassification{244.01}
\qtags{motcle=tangentes, points stationnaires}
La trajectoire $\Gamma$ d'une particule en mouvement est donnée par les équations
$$x=t^2+t^3\quad \mbox{et} \quad  y=2t^2-t^3.$$
\begin{answers}
\bad{Le point de paramètre $0$ est un point de rebroussement de seconde espèce.}
\good{La tangente à $\Gamma$ au point $(x(0),y(0))$ est la droite d'équation $y=2x$.}
\good{La tangente à $\Gamma$ au point $(x(1),y(1))$ est dirigée par le vecteur $(5,1)$.}
\bad{La tangente à $\Gamma$ au point $(x(1),y(1))$ est la droite d'équation $5x-y+3=0$.}
\end{answers}
\begin{explanations}
Le $DL_3(0)$ de $f(t)=(x(t),y(t))$ est
$$f(t)=t^2(1,2)+t^3(1,-1)+t^3\left(\varepsilon_1(t),\varepsilon_2(t)\right)\mbox{ avec }\lim _{t\to 0}\varepsilon_i(t)=0.$$
Donc le point de paramètre $0$ est un point de rebroussement de première espèce et la tangente en ce point est la droite d'équation $y=2x$. La tangente à $\Gamma$ au point de paramètre $1$ est dirigée par $f'(1)=(5,1)$.
\end{explanations}
\end{question}
}
